\documentclass{beamer}

\usepackage{amsmath,amsfonts,amssymb}
\usepackage{physics}
\usepackage{graphicx}
\usepackage{bm}

\title{The Blocking and Bootstrap Methods for Error Estimation in Monte Carlo Simulations}
\author{Morten Hjorth-Jensen}
\date{Spring 2026}

\begin{document}

\frame{\titlepage}

%------------------------------------------------

\begin{frame}{Motivation}

Monte Carlo simulations are central in many areas of physics

\begin{itemize}
\item Quantum Monte Carlo
\item Lattice gauge theory
\item Molecular dynamics
\item Statistical mechanics
\end{itemize}

Typical goal: estimate expectation values

\[
\mu = \mathbb{E}[X]
\]

from a sequence of measurements

\[
X_1, X_2, \dots, X_N
\]

Problem:

\begin{itemize}
\item Successive samples are often \textbf{correlated}
\item Naive statistical error estimates become incorrect
\end{itemize}

The blocking method provides a practical solution.
\end{frame}

%------------------------------------------------

\begin{frame}{Naive Error Estimate}

For independent samples the sample mean is

\[
\bar X = \frac{1}{N}\sum_{i=1}^{N} X_i
\]

The variance of the estimator is

\[
\mathrm{Var}(\bar X)
=
\frac{\sigma^2}{N}
\]

with

\[
\sigma^2 = \mathrm{Var}(X)
\]


\end{frame}




\begin{frame}{Naive Error Estimate, estimator}

Estimator for the error:

\[
\sigma_{\bar X}^2
=
\frac{1}{N(N-1)}
\sum_{i=1}^N (X_i-\bar X)^2
\]

This formula is only correct if the samples are \textbf{independent}.
\end{frame}


\begin{frame}{Correlated Data}

  In Markov Chain Monte Carlo we generate

  \[
  X_1 \rightarrow X_2 \rightarrow X_3 \rightarrow \dots
  \]

  via a Markov process.

  Thus

  \[
  \text{Cov}(X_i,X_j) \neq 0
  \]

  Define the autocorrelation function

  \[
  C(t)
  =
  \langle (X_i-\mu)(X_{i+t}-\mu) \rangle
  \]

\end{frame}


\begin{frame}{Correlated Data, Normalized}

  Normalized autocorrelation

  \[
  \rho(t) = \frac{C(t)}{C(0)}
  \]

  Typically

  \[
  \rho(t) \sim e^{-t/\tau}
  \]

  where $\tau$ is the autocorrelation time.
\end{frame}

%------------------------------------------------

\begin{frame}{Variance of the Mean with Correlations}

  The variance of the mean becomes

  \[
  \mathrm{Var}(\bar X)
  =
  \frac{\sigma^2}{N}
  \left(
  1
  +
  2\sum_{t=1}^{\infty}\rho(t)
  \right)
  \]

  Define the \textbf{integrated autocorrelation time}

  \[
  \tau_{\mathrm{int}}
  =
  \frac{1}{2}
  +
  \sum_{t=1}^{\infty}\rho(t)
  \]

  Then

  \[
  \mathrm{Var}(\bar X)
  =
  \frac{2\tau_{\mathrm{int}}\sigma^2}{N}
  \]

\end{frame}


\begin{frame}{Variance of the Mean with Correlations, Interpretation}

Interpretation

  \begin{itemize}
  \item Effective number of independent samples
  \end{itemize}

  \[
  N_{\mathrm{eff}}
  =
  \frac{N}{2\tau_{\mathrm{int}}}
  \]
\end{frame}

%------------------------------------------------

\begin{frame}{Idea Behind the Blocking Method}

  Instead of estimating autocorrelations explicitly:

  \begin{itemize}
  \item Combine neighboring samples into blocks
  \item Each blocking level reduces correlations
  \end{itemize}

  If blocks are large enough:

  \[
  \text{block averages} \approx \text{independent samples}
  \]

  Then the usual variance estimator becomes valid again.
\end{frame}

%------------------------------------------------

\begin{frame}{Blocking Transformation}

  Start with dataset

  \[
  X^{(0)}_i, \quad i=1,\dots,N
  \]

  Define the blocking transformation

  \[
  X^{(k+1)}_i
  =
  \frac{1}{2}
  \left(
  X^{(k)}_{2i-1}
  +
  X^{(k)}_{2i}
  \right)
  \]

\end{frame}



\begin{frame}{Blocking Transformation, Properties}

  Properties:

  \begin{itemize}
  \item Number of samples halves each iteration
  \item Correlations decrease with blocking level
  \end{itemize}

  After $k$ transformations:

  \[
  N_k = \frac{N}{2^k}
  \]
\end{frame}




\begin{frame}{Variance at Blocking Level $k$}

  Compute

  \[
  \bar X_k
  =
  \frac{1}{N_k}
  \sum_{i=1}^{N_k} X^{(k)}_i
  \]

  Variance estimate

  \[
  s_k^2
  =
  \frac{1}{N_k-1}
  \sum_{i=1}^{N_k}
  \left(
  X^{(k)}_i - \bar X_k
  \right)^2
  \]

\end{frame}



\begin{frame}{Variance at Blocking Level $k$, Error estimate}


  Error estimate

  \[
  \sigma_k^2
  =
  \frac{s_k^2}{N_k}
  \]

  As correlations disappear

  \[
  \sigma_k^2
  \rightarrow
  \mathrm{Var}(\bar X)
  \]
\end{frame}


%------------------------------------------------

\begin{frame}{Convergence of the Blocking Method}

  Key observation:

  \begin{itemize}
  \item Initially the variance estimate increases
  \item Eventually it reaches a plateau
  \end{itemize}

  Plateau value gives the correct error estimate.

  Graphically:

  \begin{itemize}
  \item Plot variance vs blocking level
  \item Choose region where curve stabilizes
  \end{itemize}

  This plateau indicates that blocks are effectively independent.
\end{frame}

%------------------------------------------------
\begin{frame}{Rigorous Derivation from the Covariance Matrix}

Let
\[
\bm{X} = (X_1,X_2,\dots,X_N)^T
\]
be a random vector with mean
\[
\mathbb{E}[\bm{X}] = \mu \bm{1},
\qquad
\bm{1} = (1,1,\dots,1)^T .
\]

The sample mean can be written as
\[
\bar X = \frac{1}{N}\bm{1}^T \bm{X}.
\]

We want the variance of $\bar X$ when the entries of $\bm{X}$ are correlated.

Define the covariance matrix
\[
\Sigma_{ij} = \mathrm{Cov}(X_i,X_j)
= \mathbb{E}\big[(X_i-\mu)(X_j-\mu)\big].
\]

Then
\[
\mathrm{Var}(\bar X)
=
\mathrm{Var}\left(\frac{1}{N}\bm{1}^T\bm{X}\right).
\]
\end{frame}

%------------------------------------------------
\begin{frame}{Variance of a Linear Combination}

For any random vector $\bm{X}$ and deterministic vector $\bm{a}$,
\[
\mathrm{Var}(\bm{a}^T\bm{X}) = \bm{a}^T \Sigma \bm{a}.
\]

Choosing
\[
\bm{a} = \frac{1}{N}\bm{1},
\]
we obtain
\[
\mathrm{Var}(\bar X)
=
\left(\frac{1}{N}\bm{1}\right)^T
\Sigma
\left(\frac{1}{N}\bm{1}\right)
=
\frac{1}{N^2}\bm{1}^T\Sigma\bm{1}.
\]

\end{frame}




%------------------------------------------------
\begin{frame}{Variance of a Linear Combination}


Writing this out explicitly,
\[
\mathrm{Var}(\bar X)
=
\frac{1}{N^2}\sum_{i=1}^N \sum_{j=1}^N \Sigma_{ij}
=
\frac{1}{N^2}\sum_{i=1}^N \sum_{j=1}^N \mathrm{Cov}(X_i,X_j).
\]

This is the exact formula for the variance of the sample mean for correlated data.
\end{frame}


%------------------------------------------------
\begin{frame}{Stationary Time Series and Toeplitz Structure}

Assume the process is weakly stationary:

\begin{itemize}
\item $\mathbb{E}[X_i]=\mu$ independent of $i$
\item $\mathrm{Cov}(X_i,X_j)$ depends only on the lag $|i-j|$
\end{itemize}

\end{frame}


%------------------------------------------------
\begin{frame}{Stationary Time Series and Toeplitz Structure}


Then
\[
\Sigma_{ij} = C(|i-j|),
\]
where
\[
C(t)=\mathrm{Cov}(X_k,X_{k+t}).
\]

Thus the covariance matrix has Toeplitz form:
\[
\Sigma =
\begin{pmatrix}
C(0) & C(1) & C(2) & \cdots \\
C(1) & C(0) & C(1) & \cdots \\
C(2) & C(1) & C(0) & \cdots \\
\vdots & \vdots & \vdots & \ddots
\end{pmatrix}.
\]

The diagonal entries are the ordinary variance:
\[
C(0)=\sigma^2.
\]
\end{frame}

%------------------------------------------------
\begin{frame}{Summing the Covariances}

Using stationarity,
\[
\mathrm{Var}(\bar X)
=
\frac{1}{N^2}\sum_{i=1}^N\sum_{j=1}^N C(|i-j|).
\]

We now reorganize the double sum according to the lag
\[
t = |i-j|.
\]

For lag $t=0$, there are $N$ terms.

For lag $t\ge 1$, there are
\[
2(N-t)
\]
terms, since both $(i,j)$ and $(j,i)$ contribute.


\end{frame}


\begin{frame}{Summing the Covariances}



Therefore
\[
\mathrm{Var}(\bar X)
=
\frac{1}{N^2}
\left[
N C(0) + 2\sum_{t=1}^{N-1}(N-t)C(t)
\right].
\]

Equivalently,
\[
\mathrm{Var}(\bar X)
=
\frac{C(0)}{N}
+
\frac{2}{N^2}\sum_{t=1}^{N-1}(N-t)C(t).
\]
\end{frame}



%------------------------------------------------
\begin{frame}{Normalized Autocorrelation Function}

Introduce the normalized autocorrelation function
\[
\rho(t)=\frac{C(t)}{C(0)}=\frac{C(t)}{\sigma^2}.
\]

Then
\[
\mathrm{Var}(\bar X)
=
\frac{\sigma^2}{N}
\left[
1+\frac{2}{N}\sum_{t=1}^{N-1}(N-t)\rho(t)
\right].
\]

Rewriting,
\[
\mathrm{Var}(\bar X)
=
\frac{\sigma^2}{N}
\left[
1+2\sum_{t=1}^{N-1}\left(1-\frac{t}{N}\right)\rho(t)
\right].
\]

This formula is exact for a stationary correlated sequence of finite length $N$.
\end{frame}

%------------------------------------------------
\begin{frame}{Large-$N$ Limit and Integrated Autocorrelation Time}

Suppose the autocorrelation decays sufficiently rapidly so that
\[
\sum_{t=1}^{\infty} |\rho(t)| < \infty.
\]

Then for large $N$,
\[
1-\frac{t}{N}\approx 1
\qquad
\text{for relevant lags } t,
\]
and hence
\[
\mathrm{Var}(\bar X)
\approx
\frac{\sigma^2}{N}
\left[
1+2\sum_{t=1}^{\infty}\rho(t)
\right].
\]


\end{frame}



\begin{frame}{Large-$N$ Limit and Integrated Autocorrelation Time}


Define the integrated autocorrelation time
\[
\tau_{\mathrm{int}}
=
\frac{1}{2}+\sum_{t=1}^{\infty}\rho(t).
\]

Then
\[
\mathrm{Var}(\bar X)
\approx
\frac{2\tau_{\mathrm{int}}\sigma^2}{N}.
\]

This shows precisely how correlations inflate the uncertainty of the mean.
\end{frame}



%------------------------------------------------
\begin{frame}{Interpretation in Terms of Effective Sample Size}

If the data were independent, we would have
\[
\mathrm{Var}(\bar X)=\frac{\sigma^2}{N}.
\]

For correlated data,
\[
\mathrm{Var}(\bar X)\approx \frac{2\tau_{\mathrm{int}}\sigma^2}{N}.
\]

Hence we define an effective number of independent samples by
\[
\frac{\sigma^2}{N_{\mathrm{eff}}}
=
\frac{2\tau_{\mathrm{int}}\sigma^2}{N}.
\]

Therefore
\[
N_{\mathrm{eff}}=\frac{N}{2\tau_{\mathrm{int}}}.
\]

So autocorrelations reduce the information content of the sample from $N$ to $N_{\mathrm{eff}}$.
\end{frame}

%------------------------------------------------
\begin{frame}{Why Blocking Works}

The blocking transformation is
\[
X_i^{(k+1)}=\frac{1}{2}\left(X_{2i-1}^{(k)}+X_{2i}^{(k)}\right).
\]

This is a linear coarse-graining of the original correlated sequence.

Its effect is to suppress short-range correlations and to reduce the off-diagonal structure of the covariance matrix.

At sufficiently high blocking level, the transformed covariance matrix becomes approximately diagonal:
\[
\Sigma^{(k)} \approx s_k^2 I.
\]

Then
\[
\mathrm{Var}(\bar X)
\approx \frac{s_k^2}{N_k},
\]
which is exactly the independent-sample formula.

Thus blocking is a practical numerical procedure for converting the full covariance structure into a reliable variance estimate.
\end{frame}

%------------------------------------------------
\begin{frame}{Mathematical Summary}

Starting from the covariance matrix,
\[
\mathrm{Var}(\bar X)=\frac{1}{N^2}\bm{1}^T\Sigma\bm{1}.
\]

For a stationary process,
\[
\Sigma_{ij}=C(|i-j|),
\]
so
\[
\mathrm{Var}(\bar X)
=
\frac{1}{N^2}
\left[
N C(0)+2\sum_{t=1}^{N-1}(N-t)C(t)
\right].
\]


\end{frame}


%------------------------------------------------
\begin{frame}{Mathematical Summary}




In terms of $\rho(t)=C(t)/C(0)$,
\[
\mathrm{Var}(\bar X)
=
\frac{\sigma^2}{N}
\left[
1+2\sum_{t=1}^{N-1}\left(1-\frac{t}{N}\right)\rho(t)
\right].
\]

For large $N$,
\[
\mathrm{Var}(\bar X)\approx \frac{2\tau_{\mathrm{int}}\sigma^2}{N}.
\]

This is the theoretical foundation behind blocking analysis.
\end{frame}



%------------------------------------------------

\begin{frame}{Blocking Algorithm}

  Given data $X_1,\dots,X_N$:

  \begin{enumerate}

  \item Compute variance estimate

    \[
    \sigma_0^2
    \]

  \item Perform blocking transformation

    \[
    X_i \rightarrow \frac{X_{2i-1}+X_{2i}}{2}
    \]

  \item Recompute variance

    \[
    \sigma_k^2
    \]

  \item Repeat until

    \[
    N_k < 2
    \]

  \item Identify plateau region in variance

  \end{enumerate}

  The plateau gives the best estimate of the statistical error.
\end{frame}

\begin{frame}{Blocking Summary}

  \begin{itemize}

  \item Monte Carlo samples are typically correlated
  \item Naive variance estimates underestimate errors
  \item Blocking reduces correlations through recursive averaging
  \item Variance converges to a plateau corresponding to correct error
  \item Simple, robust and widely used in computational physics

  \end{itemize}

  Blocking remains one of the most reliable error estimation techniques in large-scale Monte Carlo simulations.
\end{frame}

%------------------------------------------------


\begin{frame}{Basic Idea of the Bootstrap}

Suppose the unknown distribution generating the data is

\[
F(x)
\]

but we only observe samples

\[
X_1,\dots,X_N \sim F
\]

The bootstrap replaces the unknown distribution by the empirical distribution

\[
\hat{F}_N(x)
\]

defined by

\[
\hat{F}_N(x)=\frac{1}{N}\sum_{i=1}^{N}\delta(x-X_i)
\]

\end{frame}


\begin{frame}{Basic Idea of the Bootstrap, Interpretation}



Interpretation:

\begin{itemize}
\item Treat the observed dataset as the population
\item Generate new samples by resampling with replacement
\end{itemize}

This produces synthetic datasets from $\hat{F}_N$.
\end{frame}



%------------------------------------------------

\begin{frame}{Bootstrap Resampling}

Generate a bootstrap sample

\[
X_1^*, X_2^*, \dots, X_N^*
\]

by drawing $N$ values from the original data

\[
\{X_1,\dots,X_N\}
\]

with replacement.

Example:

\[
\{1.2, 0.8, 1.4, 0.9\}
\]

Possible bootstrap sample:

\[
\{0.8,1.4,1.4,1.2\}
\]

Each bootstrap dataset mimics an independent experiment.
\end{frame}

%------------------------------------------------

\begin{frame}{Bootstrap Distribution of an Estimator}

For each bootstrap sample compute

\[
\hat{\theta}^* = f(X_1^*,\dots,X_N^*)
\]

Repeat this procedure $B$ times:

\[
\hat{\theta}_1^*, \hat{\theta}_2^*, \dots, \hat{\theta}_B^*
\]

These values form the \textbf{bootstrap distribution}

\[
P_B(\hat{\theta})
\]

which approximates the sampling distribution of the estimator

\[
P(\hat{\theta})
\]

in the limit of large $N$.
\end{frame}

%------------------------------------------------

\begin{frame}{Bootstrap Estimate of the Variance}

The bootstrap estimate of the variance is

\[
\mathrm{Var}_{\text{boot}}(\hat{\theta})
=
\frac{1}{B-1}
\sum_{b=1}^{B}
(\hat{\theta}_b^*-\bar{\theta}^*)^2
\]

where

\[
\bar{\theta}^*=
\frac{1}{B}\sum_{b=1}^{B}\hat{\theta}_b^*
\]

Thus the bootstrap standard error becomes

\[
\sigma_{\text{boot}}
=
\sqrt{\mathrm{Var}_{\text{boot}}}.
\]

As $B\to\infty$, the estimate converges to the true variance.
\end{frame}

%------------------------------------------------

\begin{frame}{Consistency of the Bootstrap}

Let

\[
\hat{\theta}=f(X_1,\dots,X_N)
\]

be an estimator.

Under mild regularity conditions,

\[
\sqrt{N}(\hat{\theta}-\theta)
\]

has a limiting distribution.

Bootstrap theory shows

\[
\sqrt{N}(\hat{\theta}^*-\hat{\theta})
\]

approximates the same limiting distribution.

Thus

\[
P^*(\hat{\theta}^*-\hat{\theta})
\approx
P(\hat{\theta}-\theta).
\]

This is the theoretical justification of the bootstrap.
\end{frame}

%------------------------------------------------

\begin{frame}{Bootstrap Confidence Intervals}

Let

\[
\hat{\theta}_1^*,\dots,\hat{\theta}_B^*
\]

be sorted.

\textbf{Percentile interval}

For confidence level $1-\alpha$:

\[
[\hat{\theta}_{\alpha/2}^*, \hat{\theta}_{1-\alpha/2}^*]
\]

Example:

95\% interval

\[
[\hat{\theta}_{2.5\%}^*,\hat{\theta}_{97.5\%}^*]
\]

This method uses the empirical quantiles of the bootstrap distribution.
\end{frame}

%------------------------------------------------

\begin{frame}{Bias Estimation}

Bootstrap also estimates bias:

\[
\text{Bias} =
\mathbb{E}[\hat{\theta}] - \theta
\]

Bootstrap estimate

\[
\widehat{\text{Bias}}
=
\bar{\theta}^* - \hat{\theta}
\]

Bias-corrected estimator:

\[
\hat{\theta}_{BC}
=
\hat{\theta} - \widehat{\text{Bias}}
\]

Bootstrap therefore provides both

\begin{itemize}
\item variance estimates
\item bias corrections
\end{itemize}
\end{frame}

%------------------------------------------------

\begin{frame}{Bootstrap Algorithm}

Given data

\[
X_1,\dots,X_N
\]

\begin{enumerate}

\item Choose number of resamples $B$

\item For $b=1\dots B$

\begin{itemize}

\item Draw bootstrap sample with replacement
\[
X_1^*,\dots,X_N^*
\]

\item Compute estimator
\[
\hat{\theta}_b^*
\]

\end{itemize}

\item Compute

\[
\sigma_{\text{boot}}
\]

\end{enumerate}

Typical choice:

\[
B = 1000 - 10000
\]
\end{frame}


\begin{frame}{Applications in Physics}

Bootstrap is widely used in

\begin{itemize}

\item Lattice QCD
\item Quantum Monte Carlo
\item Molecular simulations
\item Experimental particle physics
\item Parameter estimation in model fitting

\end{itemize}

Examples of estimators:

\begin{itemize}

\item correlation functions
\item energy differences
\item nonlinear observables
\item fitted parameters
\end{itemize}
\end{frame}

%------------------------------------------------

\begin{frame}{Example: Monte Carlo Observable}

Suppose we estimate an observable

\[
\langle O \rangle
=
\frac{1}{N}\sum_i O_i
\]

But we want uncertainty for a nonlinear function

\[
R = \frac{\langle A \rangle}{\langle B \rangle}
\]

Analytic variance formulas become complicated.

Bootstrap procedure:

\begin{itemize}

\item resample pairs $(A_i,B_i)$
\item recompute ratio
\item obtain distribution of $R$

\end{itemize}

This yields robust uncertainty estimates.
\end{frame}

%------------------------------------------------

\begin{frame}{Bootstrap vs Other Resampling Methods}

Common resampling techniques:

\begin{itemize}

\item Bootstrap
\item Jackknife
\item Blocking
\end{itemize}

Bootstrap advantages:

\begin{itemize}

\item works for complex estimators
\item easy to implement
\item provides full distribution
\end{itemize}


\end{frame}


\begin{frame}{Bootstrap vs Other Resampling Methods}


Bootstrap limitation:

\begin{itemize}

\item assumes independent samples
\end{itemize}

For correlated Monte Carlo data we combine:

\begin{itemize}

\item blocking + bootstrap
\end{itemize}
\end{frame}


%------------------------------------------------

\begin{frame}{Summary Bootstrap}

Bootstrap resampling is a powerful statistical tool.

\begin{itemize}

\item Approximates sampling distributions numerically
\item Requires minimal analytical assumptions
\item Provides variance, bias, and confidence intervals
\item Widely used in computational physics

\end{itemize}

In modern simulations bootstrap analysis is a standard component of data analysis pipelines.
\end{frame}





\end{document}
